\documentclass[11pt]{article}

\usepackage[utf8]{inputenc}
\usepackage[T1]{fontenc}
\usepackage{lmodern}
\usepackage{geometry}
\usepackage{amsmath,amssymb,amsfonts,amsthm,mathtools}
\usepackage{bm}
\usepackage{enumitem}
\usepackage{microtype}
\usepackage{hyperref}
\usepackage{titlesec}
\usepackage{booktabs}
\usepackage{array}
\usepackage{setspace}
\usepackage{mathrsfs}
\usepackage{physics}

\geometry{margin=1in}
\hypersetup{colorlinks=true, linkcolor=black, urlcolor=blue, citecolor=black}
\setlist[enumerate]{topsep=0.4em,itemsep=0.35em}
\setlist[itemize]{topsep=0.3em,itemsep=0.25em}
\titleformat{\section}{\large\bfseries}{\thesection.}{0.5em}{}
\titleformat{\subsection}{\normalsize\bfseries}{\thesubsection.}{0.5em}{}
\setstretch{1.06}

\newcommand{\Aether}{\AE{}ther}
\newcommand{\Aflow}{\AE{}ther-flow}
\newcommand{\TheoryName}{\Aflow{} Interpretation of Relativity}
\newcommand{\TheoryProgram}{\Aether{} / \Aflow{} framework}
\newcommand{\CompletedMetric}{g^{\sharp}}
\newcommand{\SelectorCore}{\mathfrak C^{\mathrm{sel,split}}}
\newcommand{\ShellRestriction}{\mathfrak S^{\mathrm{sub,sh}}}
\newcommand{\LiftPackage}{\mathfrak L^{\mathrm{sub,lift}}}

\newtheorem{definition}{Definition}
\newtheorem{proposition}{Proposition}
\newtheorem{remark}{Remark}
\newtheorem{corollary}{Corollary}
\newtheorem{theorem}{Theorem}

\title{\textbf{The \TheoryName{}}\\[0.4em]
\large Primitive-Reservoir Observer-Reduction\\
\large Primitive Substrate Shell Restriction\\
\large Next Benchmark Criterion}
\author{}
\date{}

\begin{document}

\maketitle

\begin{abstract}
The current split-core routing note had fixed the honest burden below the selector-law collapse frontier: first derive or uniquely force the split zero-hidden selector-shell core \(\SelectorCore\) from still more primitive explicit substrate structure, or else prove a genuinely smaller theorem inside that split core. The new shell-restriction forcing theorem has now spent that first admissible route. The split core \(\SelectorCore\) is no longer left as an independently inserted stopping point; it is a canonical and unique output of the primitive substrate shell-restriction layer \(\ShellRestriction\).

This note records the routing consequence of that new theorem. The next honest benchmark-facing manuscript must therefore go one level below shell-restriction scope. Its first try should be to derive or uniquely force the shell-generating substrate lift package \(\LiftPackage\) from still more primitive explicit substrate structure in a way that actually changes benchmark-facing status. If that cannot be done cleanly, the legitimate fallback is a genuinely smaller theorem-level collapse, sharper necessity result, or sharper uniqueness result strictly inside the shell-restriction layer \(\ShellRestriction\) itself. Another split-core relay, selector-law restatement, combined-response restatement, ambient-package restatement, trace-core restatement, completion-core restatement, exact-shell affine nucleus restatement, or benchmark-neutral shell-restriction relay that preserves the same shell-restriction package without shrinking or forcing it further does not count next.

The benchmark boundary remains fixed throughout. The one operative metric is still exactly \(\CompletedMetric\), the ordinary matter route is unchanged, there is no second operative metric, the low-energy matter law is not enlarged, and stronger first-principles promotion language remains paused. The positive result of this note is therefore routing discipline at the new shell-restriction frontier.
\end{abstract}

\section{Introduction}

The active derivational program of \TheoryName{} is no longer at a stage where the split zero-hidden selector-shell core \(\SelectorCore\) should be treated as the default unexplained stopping point. The split-core collapse theorem had already shown that the full selector-law presentation carries no extra benchmark-facing freedom beyond \(\SelectorCore\). The new shell-restriction forcing theorem then spent the first admissible upstream option below that frontier: \(\SelectorCore\) is itself a canonical and unique output of the primitive substrate shell-restriction layer \(\ShellRestriction\).

The present note records the routing consequence of that theorem. It does not reopen the split-core mathematics, and it does not claim a completed first-principles substrate derivation of gravity. It records only which kind of next manuscript would now count as an honest benchmark-facing gain below the shell-restriction frontier.

\paragraph{Framework and claim boundary.}
\TheoryProgram{} denotes the broader \Aether{} / \Aflow{} framework built on the claims that the \Aether{} is the underlying four-dimensional substrate of reality, the \Aflow{} is its intrinsic ordered motion, observed three-dimensional space is the local experiential slice of that deeper substrate, S-time is the experienced order of change arising from matter, light, and the \Aflow{}, and observed expansion is the three-dimensional appearance of deeper four-dimensional ordered motion. Within that framework, \TheoryName{} denotes the exact-closure theory in which the effective gravitational dynamics are adopted to be exactly Einsteinian with universal matter coupling. In the active sequence, \emph{adoption} means use of that established relativistic dynamics without claiming substrate derivation, while \emph{derivation} is reserved for a first-principles recovery from explicit substrate variables.

The present note stays strictly inside that boundary. It records only which kind of next manuscript would now count as an honest benchmark-facing gain below the shell-restriction frontier.

\section{Fixed Inputs}

\begin{definition}[Shell-restriction-frontier fixed inputs]
The criterion recorded here uses the following fixed inputs.
\begin{enumerate}
    \item The benchmark exact-closure package, which fixes the public one-metric GR target of the repository.
    \item The benchmark gatekeeping layer, including the post-bridge status correction and the upstream-compression safeguard that blocks benchmark-neutral deeper-origin relays from counting as new front-line progress once the benchmark-relevant object is unchanged.
    \item The split-core collapse theorem, which proves that the full zero-hidden selector law carries no extra benchmark-facing freedom beyond the smaller split selector-shell core \(\SelectorCore\).
    \item The new shell-restriction forcing theorem, which proves that the split selector-shell core \(\SelectorCore\) is itself a canonical and unique output of the primitive substrate shell-restriction layer \(\ShellRestriction\).
    \item The preserved direct-support shell-restriction theorem, which records that \(\ShellRestriction\) is itself the structured output of a deeper shell-generating substrate lift package \(\LiftPackage\) on the same one-metric line.
\end{enumerate}
\end{definition}

\begin{remark}
The present note does not replace those manuscripts. It records the routing consequence of reading them together under the active safeguard against recursive depth drift.
\end{remark}

\section{Why Shell Restriction Is Now the Frontier}

\begin{definition}[Benchmark-neutral shell-restriction relay]
A \emph{benchmark-neutral shell-restriction relay} is a manuscript that preserves the same benchmark one-metric output, the same ordinary matter route, the same split selector-shell core \(\SelectorCore\), the same selector-law consequence, the same combined-response consequence, the same ambient-package consequence, the same trace-core consequence, the same completion-core consequence, and the same claim boundary while merely restating \(\SelectorCore\), the selector law, the combined response law, the ambient package, the trace core, the completion core, the exact-shell affine nucleus, or the same shell-restriction package \(\ShellRestriction\) without shrinking or forcing \(\ShellRestriction\) further.
\end{definition}

\begin{proposition}[The shell-restriction forcing theorem is the live frontier]
At the current stage of the primitive-reservoir observer-reduction line, the theorem forcing the split selector-shell core \(\SelectorCore\) from the primitive substrate shell-restriction layer \(\ShellRestriction\) is the live benchmark-facing frontier.
\end{proposition}

\begin{proof}
That theorem changes benchmark-facing status at current scope because it removes \(\SelectorCore\) itself as the live internal stopping point. The split core is no longer merely the smaller datum singled out inside the selector law; it is now forced by the deeper shell-restriction layer. By contrast, a manuscript that preserves the same shell-restriction package and the same split-core consequence while merely restating \(\SelectorCore\) or another already-forced larger or smaller bridge object does not alter benchmark-facing status unless it adds a comparably sharp gain. Therefore the live frontier is the shell-restriction forcing theorem rather than the earlier split-core collapse theorem or a later same shell-restriction relay.
\end{proof}

\begin{corollary}[Status of the split core as a next-step target]
Under the current routing safeguard, the split zero-hidden selector-shell core \(\SelectorCore\) is no longer itself the default next benchmark-facing target.
\end{corollary}

\begin{proof}
The shell-restriction forcing theorem already proves that \(\SelectorCore\) is a canonical and unique output of \(\ShellRestriction\). Once that is on record, another manuscript that spends the move on \(\SelectorCore\) again does not remove any further benchmark-relevant freedom. The open burden has therefore moved below split-core scope.
\end{proof}

\section{The Next Benchmark Criterion}

\begin{definition}[Admissible next benchmark-facing shell-restriction manuscript]
At the current shell-restriction frontier, a manuscript counts as an admissible next benchmark-facing move only if it respects the following burden order:
\begin{enumerate}
    \item its primary attempt is to derive or uniquely force the shell-generating substrate lift package \(\LiftPackage\) from still more primitive explicit substrate structure in a way that does more than merely relocate the unexplained layer deeper;
    \item if that cannot be done cleanly at current scope, its admissible fallback gain is to prove a genuinely smaller theorem-level collapse, sharper necessity result, or sharper uniqueness result strictly inside the shell-restriction layer \(\ShellRestriction\);
    \item in either case it preserves the same benchmark one-metric package, the same ordinary matter route, and the same claim boundary, and it introduces neither a second operative metric nor an enlarged low-energy matter law and licenses no stronger first-principles promotion language.
\end{enumerate}
\end{definition}

\begin{theorem}[Next honest benchmark-facing task at shell-restriction scope]
The next honest benchmark-facing manuscript at the current shell-restriction frontier should therefore go one level below that frontier: first try to derive or uniquely force the shell-generating substrate lift package \(\LiftPackage\) from still more primitive explicit substrate structure in a way that actually changes benchmark-facing status. If that cannot be done cleanly, the admissible fallback is to prove a genuinely smaller collapse, sharper necessity result, or sharper uniqueness result strictly inside \(\ShellRestriction\). A manuscript that only restates the split selector-shell core \(\SelectorCore\), only restates the selector law, only restates the combined response law, only restates the ambient package, only restates the trace core, only restates the full completion-core presentation, only restates the exact-shell affine nucleus, or only repeats the same shell-restriction consequence while preserving \(\ShellRestriction\) without shrinking or forcing it further does not satisfy the criterion.
\end{theorem}

\begin{proof}
By Proposition 1, the live frontier already lies at the theorem that forces the split core from shell restriction. Therefore a second manuscript below it must improve on that status rather than merely accompany it. A deeper-origin manuscript can do so only if it changes benchmark-facing status by more than depth alone, and the first clean way to do that is to force the shell-generating substrate lift package that already sits immediately upstream in the direct-support record. If that first option is not yet available cleanly, the routing safeguard allows the fallback of a smaller internal collapse, necessity result, or uniqueness result inside \(\ShellRestriction\). If a manuscript simply reproduces the same shell-restriction consequence through another relay while preserving the same shell-restriction package without shrinking or forcing it further, the routing rule classifies it as benchmark neutral.
\end{proof}

\section{What the Next Move Should Not Be}

\begin{proposition}[Screened next moves]
At the current shell-restriction frontier, none of the following is the default next benchmark-facing move:
\begin{enumerate}
    \item another restatement of the split zero-hidden selector-shell core \(\SelectorCore\);
    \item another restatement of the full zero-hidden selector law;
    \item another restatement of the combined response substrate law;
    \item another restatement of the completion-balance ambient package;
    \item another restatement of the completion-state shell-trace core;
    \item another restatement of the full completion-core presentation;
    \item another restatement of the exact-shell affine nucleus;
    \item another shell-restriction relay that preserves the same \(\ShellRestriction\) without shrinking or forcing it further;
    \item a manuscript that introduces a second operative metric, enlarges the ordinary matter law, or uses stronger first-principles promotion language.
\end{enumerate}
\end{proposition}

\begin{proof}
Items (1)--(8) fail the preceding criterion because they do not directly address the current shell-restriction burden. They revisit either already-forced larger bridge objects, already-forced smaller bridge objects, or the same shell-restriction stopping point rather than removing the remaining shell-restriction freedom or proving a still smaller internal datum inside \(\ShellRestriction\). Item (9) violates the fixed claim boundary of the benchmark package and therefore cannot count as a disciplined next step on the active line. The benchmark package remains the exact one-metric GR package already fixed by adoption, so any admissible next manuscript must preserve that boundary while keeping stronger first-principles promotion language paused.
\end{proof}

\begin{remark}
This screening statement is not a denial that the selector-law collapse theorem, the split-core forcing theorem, or the preserved shell-restriction support theorem matter historically. It is a routing statement: they are not the honest way to spend the next benchmark-facing move from the current shell-restriction frontier.
\end{remark}

\section{Conclusion}

The positive result of this note is routing discipline at the current shell-restriction frontier. The shell-restriction forcing theorem remains the live benchmark-facing gain because it already removes the split selector-shell core as the internal stopping point on the active line. The split core therefore no longer sets the honest next burden, and neither do the already-forced selector law, combined response law, ambient package, trace core, completion-core presentation, or exact-shell affine nucleus.

The scope remains narrow and honest. The benchmark package is unchanged: the one operative metric is still exactly \(\CompletedMetric\), the ordinary matter route is unchanged, there is no second operative metric, and the low-energy matter law is not enlarged. Nothing here upgrades adoption to derivation or licenses stronger first-principles promotion language yet.

The next open burden is therefore clear. The next benchmark-facing manuscript should go one level below the current shell-restriction frontier: first try to derive or uniquely force the shell-generating substrate lift package \(\LiftPackage\) from still more primitive explicit substrate structure in a way that actually changes benchmark-facing status while preserving the same benchmark one-metric package and claim boundary. If that cannot be done cleanly, the admissible fallback is a genuinely smaller collapse, sharper necessity result, or sharper uniqueness result strictly inside the shell-restriction layer \(\ShellRestriction\). Another split-core relay, selector-law restatement, combined-response restatement, ambient-package restatement, trace-core restatement, completion-core restatement, exact-shell affine nucleus restatement, or another shell-restriction relay that preserves the same shell-restriction package without shrinking or forcing it further is not the clean next manuscript target at current scope.

\end{document}
